\newpage
\section{Python}
\index{Python}
\nocite{py3:book}


\begin{center}
	\includegraphics[scale=0.5]{pictures/python/logo.pdf}
\end{center}

Python je objektově orientovaný a interpretovaný programovací jazyk s dynamickým a silným typováním. Python je charakteristický pro své vyjadřování struktury kódu pomocí odsazování, jehož dodržování je povinné. To vede k čitelnějšímu a přehlednějšímu kódu. 

První verze jazyka byla uvolněna v roce 1991. Python navrhl Guido van Rossum, který byl inspirován jazyky jako C++ či Perl. Jedná se o \textbf{open source} projekt dostupný pod licencí Python Software Foundation License, která je kompatibilní s GPL licencí. Rozdíl je v tom, že u Python Software Foundation License můžeme měnit kód bez nutnosti zveřejnit změny jako open source [viz \footnote{http://docs.python.org/license.html}]. Aktuální stabilní verze jsou 2.7.3 a 3.2.3 pro verzi Python 3.0, která byla uvolněna v roce 2008. 


% noveho v python3 http://docs.python.org/release/3.0.1/whatsnew/3.0.html

V dnešní době se s Pythonem můžeme setkat téměř všude. Pro jeho jednoduchost při psaní kódu je velmi populární a široce rozšířený. Jeho velká výhoda je tedy velmi čitelný kód a rychlost psaní. Jazyk je jednoduše přenositelný (není nutná kompilace), je multiplatformní a má kvalitní dokumentaci. 

V mnoha projektech existuje skriptovací rozhraní pro Python. Jako příklad uveďme QGIS, GIMP, Inkscape, Scribus, LibreOffice, Blender nebo ArcGIS. Můžeme jej najít v projektech jako Maya, OpenShot Video Editor, Wammu, DopBox, MapServer či Gajim. Používá se pro psaní grafického rozhraní. Existují verze grafických knihoven Qt a GTK pro Python - PyQt, resp. PyGTK. Využívá se jako skriptovací jazyk pro psaní webových aplikací. V síti Internet se s ním můžeme také setkat v podobě nástrojů jako Django, Zope či Pylons. Byly napsány knihovny pro vědecké výpočty - NumPy, SciPy, Matplotlib.

%Řekl bych, že odinstalovat Python z Vašeho počítače nebude zrovna ten nejlepší nápad. Je totiž prakticky všude.

\subsection*{PyQGIS}
Quantum GIS nabízí podporu QGIS API pro jazyk Python. Jedná se o verzi PyQGIS. PyQGIS můžeme používat přímo v QGISu přes příkazovou řádku, psát zásuvné moduly či využít QGIS API pro náš vlastní program. 

Existuje velmi dobře zpracovaná tzv. kuchařka jak psát v PyQGIS \cite{pyqgis:www}, tudíž nepovažuji za nutné se mnoho rozepisovat. Zmíním jen pár informací, které považuji za základní. Rastrové a vektorové vrstvy jsou reprezentovány třídami \textit{QgsRasterLayer} a \textit{QgsVectorLayer}. Třída \textit{QgsMapLayerRegistry} slouží k načítání, správu a mazání geografických vrstev z QGISu. Třída se nachází v \textbf{qgis.core} knihovně. Pomocí příkazu \textbf{QgsMapLayer.instance().mapLayers()} získáme všechny vrstvy momentálně načtené v QGIS. Aktivní vrstvu získáme pomocí příkazu \textbf{qgis.utils.iface.activeLayer()}. Vrstvu přidáme do QGISu pomocí \textbf{QgsMapLayer.instance().addMapLayer(layer)} a odebereme \textbf{QgsMapLayer.instance().removeMapLayer(layer)}. 

Vektorovou vrstvu vytvoříme takto \textbf{QgsVectorLayer}(cesta\_k\_souboru, jmeno\_vrstvy, knihovna).

Například:

\begin{lstlisting}
	vector = QgsVectorLayer("~/geodata/contour.shp", "ContourLines", "ogr")
\end{lstlisting}

Rastrovou vrstvu vytvoříme takto \textbf{QgsRasterLayer}(cesta\_k\_souboru, jmeno\_vrstvy).

Například:

\begin{lstlisting}
	raster = QgsRasterLayer("~/geodata/elevation.tiff", "Elevation")
\end{lstlisting}

Z mapových vrstev můžeme získat například souřadnicový systém pomocí metody \textit{crs}(), zdrojový soubor metodou \textit{source}(), jménou \textit{name}() či rozsah \textit{exent}().\\

Kompletní dokumentace QGIS API je dostupná zde \cite{qgis_api:www}.
%QGIS používá verzi 2.7.1.

%{\color{red}napsat vie o PyQGIS}






